\documentclass[11pt,landscape]{article}
\usepackage[french]{babel}
\usepackage[utf8]{inputenc}
\usepackage[T1]{fontenc}
\usepackage{xspace}
\usepackage{array}
\usepackage{color}
\usepackage{listings}
\usepackage{fancyvrb} %verbatim
\usepackage{tabularx}
\usepackage{eurosym}
\usepackage{mathrsfs}
\usepackage{stmaryrd}
\usepackage{amssymb}
\usepackage[?]{amsmath}
\usepackage{epsfig}
\usepackage{float}
\usepackage{wrapfig}
\usepackage{pifont}
\usepackage{enumerate}
\usepackage[np]{numprint} %\np{12345} donne 12 345 
\usepackage{multicol}
\newcommand{\second}{2\up{d}\xspace}
\newcommand{\seconde}{2\up{de}\xspace}
\newcommand{\R}{\mathbb R}
\newcommand{\Rp}{\R_+}
\newcommand{\Rpe}{\R_+^*}
\newcommand{\Rm}{\R_-}
\newcommand{\Rme}{\R_-^*}
\newcommand{\N}{\mathbb N}
\newcommand{\D}{\mathbb D}
\newcommand{\Q}{\mathbb Q}
\newcommand{\Z}{\mathbb Z}
\newcommand{\C}{\mathbb C}
\newcommand{\grs}{\mathfrak S}
\newcommand{\IN}[1]{\llbracket 1,#1\rrbracket}
\newcommand{\card}{\text{Card}\,}
\usepackage{mathrsfs}
\newcommand{\parties}{\mathscr P}
\renewcommand{\epsilon}{\varepsilon}
\newcommand{\rmd}{\text{d}}
\newcommand{\diff}{\mathrm D}
\newcommand{\Id}{{\rm Id}}
\newcommand{\e}{{\rm e}}
\newcommand{\I}{{\rm i}}
\newcommand{\J}{{\rm j}}
\newcommand{\ro}{\circ}
\newcommand{\exu}{\exists\,!\,}
\newcommand{\telq}{\,\, \mid \,\,}
\newcommand{\para}{\raisebox{0.1em}{\text{\footnotesize /\hspace{-0.1em}/}}}   
\newcommand{\vect}[1]{\overrightarrow{#1}}
\newcommand{\scal}[2]{\left(\, #1 \mid #2 \, \right)}
\newcommand{\ortho}[1]{{#1}^\perp}
\newcommand{\veci}{\vec{\text{\it \i}}}
\newcommand{\vecj}{\vec{\text{\it \j}}}
\newcommand{\rep}{$(O;\veci,\vecj,\vec{k})$\xspace}
\newcommand{\Oijk}{$(O, \veci,\vecj,\vec{k})$\xspace}
\newcommand{\rond}{repère orthonormal direct}
\newcommand{\bond}{base orthonormale directe}
\newcommand{\eq}{\Longleftrightarrow}
\newcommand{\implique}{\Longrightarrow}
\newcommand{\noneq}{\ \ \ /\hspace{-1.45em}\eq}
\newcommand{\tend}{\longrightarrow}
\newcommand{\egx}[2]{\underset{#1 \tend #2}=}
\newcommand{\asso}{\longmapsto}
\newcommand{\vers}{\longrightarrow}
\newcommand{\eqn}{~\underset{n \rightarrow \infty}{\sim}~}
\newcommand{\eqx}[2]{~\underset{#1 \rightarrow #2}{\sim}~}
\newcommand{\egn}{~\underset{n \rightarrow \infty}{=}~}
\renewcommand{\descriptionlabel}{\hspace{\labelsep}$\bullet$}
\renewcommand{\bar}{\overline}
\DeclareMathOperator{\ash}{{Argsh}}
\DeclareMathOperator{\cotan}{{cotan}}
\DeclareMathOperator{\ach}{{Argch}}
\DeclareMathOperator{\ath}{{Argth}}
\DeclareMathOperator{\sh}{{sh}}
\DeclareMathOperator{\ch}{{ch}}
%\DeclareMathOperator{\tanh}{{th}}
\DeclareMathOperator{\Mat}{{Mat}}
\DeclareMathOperator{\Vect}{{Vect}}
\DeclareMathOperator{\trace}{{tr}}
\newcommand{\tr}{{}^{\mathrm t}}
\newcommand{\divi}{~\big|~}
\newcommand{\ndivi}{~\not{\big|}~}
\newcommand{\et}{\wedge}
\newcommand{\ou}{\vee}
%\DeclareMathOperator{\trace}{{trace}}
\renewcommand{\det}{\operatorname{\text{dét}}}
%\DeclareMathOperator{\det}{{dét}}
\DeclareMathOperator{\grad}{{grad}}
%\DeclareMathOperator{\th}{th}
%\renewcommand{\tanh}{\mathop{\mathrm{th}}{}}
\renewcommand{\arcsin}{{\mathop{\mathrm{Arcsin}}}}
\renewcommand{\arccos}{{\mathop{\mathrm{Arccos}}}}
\renewcommand{\arctan}{{\mathop{\mathrm{Arctan}}}}
%\DeclareMathOperator{\arctan}{{Arctan}}
\renewcommand{\tanh}{{\mathop{\mathrm{th}}}}
\newcommand{\pgcd}{\mathop{\mathrm{pgcd}}}
\newcommand{\ppcm}{\mathop{\mathrm{ppcm}}}
\newcommand{\fonc}[4]{\left\{\begin{tabular}{ccc}$#1$ & $\vers$ & $#2$ \\
$#3$ & $\asso$ & $#4$ \end{tabular}\right.}
\renewcommand{\geq}{\geqslant}
\renewcommand{\leq}{\leqslant}
\renewcommand{\Re}{\text{\rm Re}}
\renewcommand{\Im}{\text{\rm Im}}
\renewcommand{\ker}{\mathop{\mathrm{Ker}}}
\newcommand{\Lin}{\mathcal L}
\newcommand{\GO}{\mathcal O}
\newcommand{\GSO}{\mathcal{SO}}
\newcommand{\GL}{\mathcal{GL}}
\renewcommand{\emptyset}{\varnothing}
\newcommand{\arc}[1]{\overset{\frown}{#1}}
\newcommand{\rg}{\mathop{\mathrm{rg}}}
\newcommand{\ds}{\displaystyle}
\newcommand{\co}[3]{\begin{pmatrix}#1 \\ #2 \\ #3\end{pmatrix}}
\newcommand{\demi}{\frac 1 2}

\newcommand{\exod}[1]{\addtocounter{cexo}{1} \vskip 1em 
\noindent{\bf \large Exercice \thecexo.\ #1}}
\newcounter{cact}
\renewcommand{\thecact}{\arabic{cact}}
\newcommand{\act}{\addtocounter{cact}{1} \vskip 1em \noindent{\bf Activité \thecact.\ }}
%\renewcommand{\labelitemi}{\m@th$\star$}
\newcommand{\limi}[2]{\underset{#1 \rightarrow #2}\lim}    
\newcommand{\nterme}[1]{{\em #1}}
\newcommand{\dif}[1]{#1}
\renewcommand{\FrenchLabelItem}{\textbullet} %bullet dans les items
%\renewcommand{\labelitemi}{\textbullet}
%\renewcommand{\labelitemii}{\textbullet}
%\renewcommand{\labelitemiii}{\textbullet}
%\renewcommand{\theenumi}{\alph{enumi}}
%\setenumerate[0]{label=\alph}
\usepackage{xr}
\usepackage{hyperref}
\usepackage{textcomp} %pour les ' dans lstlisting
\usepackage{listings}
%------------------ info
\lstset{
  language=Python,
  commentstyle=\color{red},
  upquote=true,
  literate={é}{{\'e}}1
           {ê}{{\^e}}1
           {Ã }{{\`a}}1
           {è}{{\`e}}1
           {î}{{\^i}}1
           {ï}{{\"i}}1
           {ë}{{\"e}}1
           {ù}{{\`u}}1,  
  columns=flexible,
  basicstyle=\ttfamily,
  keywordstyle=\color{blue},
  %identifierstyle=\color{green},
  stringstyle=\ttfamily,
  showstringspaces=false,
  %numbers=left, 
  %numberstyle=\tiny, 
  %stepnumber=1, 
  %numbersep=5pt,
  frame=L
}

 
\ifdefined\HCode
  \def\pgfsysdriver{pgfsys-dvisvgm4ht.def}
\fi 
\usepackage{pgf,tikz,pgfplots}
\usepackage{tkz-tab,tkz-fct}
\pgfplotsset{compat=1.16}
%\usepackage{mathrsfs}
\usetikzlibrary{arrows}

%-----------------------
\usepackage{fancybox}
\usepackage{fancyhdr}


\definecolor{bleusad}{HTML}{0056B3}
\newcommand{\titreFiche}[3]{\centerline{\Ovalbox{\parbox[t]{0.7\textwidth}{
\vspace{-0.5em}\begin{center} \bf \large #1
\end{center}
\vspace{-0.5em}}}}

  \bigskip
\fancyfoot[L]{\includegraphics[height=1em]{/var/www/sacado/static/img/sacadoA1.png}\textcolor{bleusad}{\textsf{\ \raisebox{2pt}{ {{\footnotesize https://sacado.xyz > #3}}}}}}
\fancypagestyle{Page1}{\fancyhead{}\renewcommand{\headrulewidth}{0pt}}
\thispagestyle{Page1}
\fancyhead[L]{#1}
\fancyfoot[R]{\textsc{#2}}
\renewcommand{\footrulewidth}{0.4pt}
\pagestyle{fancy}
}

\newcounter{cexo}
\renewcommand{\thecexo}{\arabic{cexo}}
\newcommand{\exo}{\refstepcounter{cexo} \vskip 1.5em \noindent{\bf Exercice \thecexo.\ }}



\definecolor{bleu2}{rgb}{0,0.58,0.79}
\definecolor{violet}{rgb}{0.79,0.66,0.98}

\newcommand{\exercice}[1]{ % Exercice directement dans la page
\vspace{0.2cm}
\tikz\node[rounded corners=2pt, fill=bleu2]{\color{white}\textbf{#1}};
\vspace{0.1cm}
}

\newcommand{\sacado}[1]{ % Exercice directement dans la page
\vspace{0.2cm}
\tikz\node[rounded corners=2pt, fill=violet]{\color{white}\textbf{#1}};
}


\newcommand{\competence}[1]{
  \hfill\textcolor{bleusad}{#1}
}
\newcommand{\competences}[1]{
  {\bf Compétences ciblées : }
  \begin{enumerate}
   \item   #1
  \end{enumerate}

  \medskip
}

\newcommand{\savoir}[1]{
  {\bf Autre savoir-faire ciblé : } 
  #1
  \medskip
}

\newcommand{\savoirs}[1]{
  %  {\bf Autres savoir-faire ciblés : }
 \vspace{-1em}
  \textcolor{bleusad}{\it \begin{itemize}
     #1
  \end{itemize}
  \centerline{\rule[0.5em]{5cm}{0.5pt}}}
  \medskip
}


\newcommand{\entreexo}{ % entre les exos}
  \medskip
  \centerline{\hrule{5cm}{1pt}}
  \medskip
}

\setlength{\parindent}{0cm}
\usepackage[left=1cm,right=1cm,top=0.5cm,bottom=1.5cm]{geometry}
\usepackage{multicol}


\newenvironment{GeneriqueT}[2]{%
\medskip \begin{tcolorbox}[widget,colback=white ,colframe=black,
title= \stepcounter{cpt}  #1  \thecpt \;: #2]}
{%
\end{tcolorbox}\par}

\usepackage{multido}
\newcommand{\point}[1]{\vspace{0.1cm}\multido{}{#1}{ \dotfill \medskip \endgraf}}
\newcommand{\ligne}[1]{\vspace{0.1cm}\multido{}{#1}{ {\color{cqcqcq}\hrulefill} \medskip \endgraf}}